\IfFileExists{stacks-project.cls}{%
\documentclass{stacks-project}
}{%
\documentclass{amsart}
}

% For dealing with references we use the comment environment
\usepackage{verbatim}
\newenvironment{reference}{\comment}{\endcomment}
%\newenvironment{reference}{}{}
\newenvironment{slogan}{\comment}{\endcomment}
\newenvironment{history}{\comment}{\endcomment}

% For commutative diagrams we use Xy-pic
\usepackage[all]{xy}

% We use 2cell for 2-commutative diagrams.
\xyoption{2cell}
\UseAllTwocells

% We use multicol for the list of chapters between chapters
\usepackage{multicol}

% This is generally recommended for better output
\usepackage{lmodern}
\usepackage[T1]{fontenc}

% For cross-file-references
\usepackage{xr-hyper}

% Package for hypertext links:
\usepackage{hyperref}

% For any local file, say "hello.tex" you want to link to please
% use \externaldocument[hello-]{hello}
\externaldocument[introduction-]{introduction}
\externaldocument[conventions-]{conventions}
\externaldocument[sets-]{sets}
\externaldocument[categories-]{categories}
\externaldocument[topology-]{topology}
\externaldocument[sheaves-]{sheaves}
\externaldocument[sites-]{sites}
\externaldocument[stacks-]{stacks}
\externaldocument[fields-]{fields}
\externaldocument[algebra-]{algebra}
\externaldocument[brauer-]{brauer}
\externaldocument[homology-]{homology}
\externaldocument[derived-]{derived}
\externaldocument[simplicial-]{simplicial}
\externaldocument[more-algebra-]{more-algebra}
\externaldocument[smoothing-]{smoothing}
\externaldocument[modules-]{modules}
\externaldocument[sites-modules-]{sites-modules}
\externaldocument[injectives-]{injectives}
\externaldocument[cohomology-]{cohomology}
\externaldocument[sites-cohomology-]{sites-cohomology}
\externaldocument[dga-]{dga}
\externaldocument[dpa-]{dpa}
\externaldocument[sdga-]{sdga}
\externaldocument[hypercovering-]{hypercovering}
\externaldocument[schemes-]{schemes}
\externaldocument[constructions-]{constructions}
\externaldocument[properties-]{properties}
\externaldocument[morphisms-]{morphisms}
\externaldocument[coherent-]{coherent}
\externaldocument[divisors-]{divisors}
\externaldocument[limits-]{limits}
\externaldocument[varieties-]{varieties}
\externaldocument[topologies-]{topologies}
\externaldocument[descent-]{descent}
\externaldocument[perfect-]{perfect}
\externaldocument[more-morphisms-]{more-morphisms}
\externaldocument[flat-]{flat}
\externaldocument[groupoids-]{groupoids}
\externaldocument[more-groupoids-]{more-groupoids}
\externaldocument[etale-]{etale}
\externaldocument[chow-]{chow}
\externaldocument[intersection-]{intersection}
\externaldocument[pic-]{pic}
\externaldocument[weil-]{weil}
\externaldocument[adequate-]{adequate}
\externaldocument[dualizing-]{dualizing}
\externaldocument[duality-]{duality}
\externaldocument[discriminant-]{discriminant}
\externaldocument[derham-]{derham}
\externaldocument[local-cohomology-]{local-cohomology}
\externaldocument[algebraization-]{algebraization}
\externaldocument[curves-]{curves}
\externaldocument[resolve-]{resolve}
\externaldocument[models-]{models}
\externaldocument[functors-]{functors}
\externaldocument[equiv-]{equiv}
\externaldocument[pione-]{pione}
\externaldocument[etale-cohomology-]{etale-cohomology}
\externaldocument[proetale-]{proetale}
\externaldocument[relative-cycles-]{relative-cycles}
\externaldocument[more-etale-]{more-etale}
\externaldocument[trace-]{trace}
\externaldocument[crystalline-]{crystalline}
\externaldocument[spaces-]{spaces}
\externaldocument[spaces-properties-]{spaces-properties}
\externaldocument[spaces-morphisms-]{spaces-morphisms}
\externaldocument[decent-spaces-]{decent-spaces}
\externaldocument[spaces-cohomology-]{spaces-cohomology}
\externaldocument[spaces-limits-]{spaces-limits}
\externaldocument[spaces-divisors-]{spaces-divisors}
\externaldocument[spaces-over-fields-]{spaces-over-fields}
\externaldocument[spaces-topologies-]{spaces-topologies}
\externaldocument[spaces-descent-]{spaces-descent}
\externaldocument[spaces-perfect-]{spaces-perfect}
\externaldocument[spaces-more-morphisms-]{spaces-more-morphisms}
\externaldocument[spaces-flat-]{spaces-flat}
\externaldocument[spaces-groupoids-]{spaces-groupoids}
\externaldocument[spaces-more-groupoids-]{spaces-more-groupoids}
\externaldocument[bootstrap-]{bootstrap}
\externaldocument[spaces-pushouts-]{spaces-pushouts}
\externaldocument[spaces-chow-]{spaces-chow}
\externaldocument[groupoids-quotients-]{groupoids-quotients}
\externaldocument[spaces-more-cohomology-]{spaces-more-cohomology}
\externaldocument[spaces-simplicial-]{spaces-simplicial}
\externaldocument[spaces-duality-]{spaces-duality}
\externaldocument[formal-spaces-]{formal-spaces}
\externaldocument[restricted-]{restricted}
\externaldocument[spaces-resolve-]{spaces-resolve}
\externaldocument[formal-defos-]{formal-defos}
\externaldocument[defos-]{defos}
\externaldocument[cotangent-]{cotangent}
\externaldocument[examples-defos-]{examples-defos}
\externaldocument[algebraic-]{algebraic}
\externaldocument[examples-stacks-]{examples-stacks}
\externaldocument[stacks-sheaves-]{stacks-sheaves}
\externaldocument[criteria-]{criteria}
\externaldocument[artin-]{artin}
\externaldocument[quot-]{quot}
\externaldocument[stacks-properties-]{stacks-properties}
\externaldocument[stacks-morphisms-]{stacks-morphisms}
\externaldocument[stacks-limits-]{stacks-limits}
\externaldocument[stacks-cohomology-]{stacks-cohomology}
\externaldocument[stacks-perfect-]{stacks-perfect}
\externaldocument[stacks-introduction-]{stacks-introduction}
\externaldocument[stacks-more-morphisms-]{stacks-more-morphisms}
\externaldocument[stacks-geometry-]{stacks-geometry}
\externaldocument[moduli-]{moduli}
\externaldocument[moduli-curves-]{moduli-curves}
\externaldocument[examples-]{examples}
\externaldocument[exercises-]{exercises}
\externaldocument[guide-]{guide}
\externaldocument[desirables-]{desirables}
\externaldocument[coding-]{coding}
\externaldocument[obsolete-]{obsolete}
\externaldocument[fdl-]{fdl}
\externaldocument[index-]{index}

% Theorem environments.
%
\theoremstyle{plain}
\newtheorem{theorem}[subsection]{Theorem}
\newtheorem{proposition}[subsection]{Proposition}
\newtheorem{lemma}[subsection]{Lemma}

\theoremstyle{definition}
\newtheorem{definition}[subsection]{Definition}
\newtheorem{example}[subsection]{Example}
\newtheorem{exercise}[subsection]{Exercise}
\newtheorem{situation}[subsection]{Situation}

\theoremstyle{remark}
\newtheorem{remark}[subsection]{Remark}
\newtheorem{remarks}[subsection]{Remarks}

\numberwithin{equation}{subsection}

% Macros
%
\def\lim{\mathop{\mathrm{lim}}\nolimits}
\def\colim{\mathop{\mathrm{colim}}\nolimits}
\def\Spec{\mathop{\mathrm{Spec}}}
\def\Hom{\mathop{\mathrm{Hom}}\nolimits}
\def\Ext{\mathop{\mathrm{Ext}}\nolimits}
\def\SheafHom{\mathop{\mathcal{H}\!\mathit{om}}\nolimits}
\def\SheafExt{\mathop{\mathcal{E}\!\mathit{xt}}\nolimits}
\def\Sch{\mathit{Sch}}
\def\Mor{\mathop{\mathrm{Mor}}\nolimits}
\def\Ob{\mathop{\mathrm{Ob}}\nolimits}
\def\Sh{\mathop{\mathit{Sh}}\nolimits}
\def\NL{\mathop{N\!L}\nolimits}
\def\CH{\mathop{\mathrm{CH}}\nolimits}
\def\proetale{{pro\text{-}\acute{e}tale}}
\def\etale{{\acute{e}tale}}
\def\QCoh{\mathit{QCoh}}
\def\Ker{\mathop{\mathrm{Ker}}}
\def\Im{\mathop{\mathrm{Im}}}
\def\Coker{\mathop{\mathrm{Coker}}}
\def\Coim{\mathop{\mathrm{Coim}}}

% Boxtimes
%
\DeclareMathSymbol{\boxtimes}{\mathbin}{AMSa}{"02}

%
% Macros for moduli stacks/spaces
%
\def\QCohstack{\mathcal{QC}\!\mathit{oh}}
\def\Cohstack{\mathcal{C}\!\mathit{oh}}
\def\Spacesstack{\mathcal{S}\!\mathit{paces}}
\def\Quotfunctor{\mathrm{Quot}}
\def\Hilbfunctor{\mathrm{Hilb}}
\def\Curvesstack{\mathcal{C}\!\mathit{urves}}
\def\Polarizedstack{\mathcal{P}\!\mathit{olarized}}
\def\Complexesstack{\mathcal{C}\!\mathit{omplexes}}
% \Pic is the operator that assigns to X its picard group, usage \Pic(X)
% \Picardstack_{X/B} denotes the Picard stack of X over B
% \Picardfunctor_{X/B} denotes the Picard functor of X over B
\def\Pic{\mathop{\mathrm{Pic}}\nolimits}
\def\Picardstack{\mathcal{P}\!\mathit{ic}}
\def\Picardfunctor{\mathrm{Pic}}
\def\Deformationcategory{\mathcal{D}\!\mathit{ef}}
\usepackage{cleveref}
\crefname{originallemma}{original lemma}{original lemmas}
\Crefname{originallemma}{Original lemma}{Original lemmas}
\makeatletter
\newlabel{lemma-bounded-fibres}{{03JX}{source}{}{sourcepointer-1}{}}
\newlabel{lemma-bounded-fibres@cref}{{[originallemma][0][]03JX}{[1][0][]source}{}{}{}}
\makeatother
\crefname{originalreference}{original reference}{original references}
\Crefname{originalreference}{Original reference}{Original references}
\makeatletter
\newlabel{morphisms-lemma-composition-etale}{{02GN}{source}{}{sourcepointer-2}{}}
\newlabel{morphisms-lemma-composition-etale@cref}{{[originalreference][0][]02GN}{[1][0][]source}{}{}{}}
\makeatother
\crefname{originalreference}{original reference}{original references}
\Crefname{originalreference}{Original reference}{Original references}
\makeatletter
\newlabel{morphisms-lemma-open-immersion-etale}{{02GP}{source}{}{sourcepointer-3}{}}
\newlabel{morphisms-lemma-open-immersion-etale@cref}{{[originalreference][0][]02GP}{[1][0][]source}{}{}{}}
\makeatother
\crefname{originalreference}{original reference}{original references}
\Crefname{originalreference}{Original reference}{Original references}
\makeatletter
\newlabel{spaces-lemma-composition-representable-transformations-property}{{02WK}{source}{}{sourcepointer-4}{}}
\newlabel{spaces-lemma-composition-representable-transformations-property@cref}{{[originalreference][0][]02WK}{[1][0][]source}{}{}{}}
\makeatother
\crefname{originalreference}{original reference}{original references}
\Crefname{originalreference}{Original reference}{Original references}
\makeatletter
\newlabel{spaces-lemma-morphism-schemes-gives-representable-transformation-property}{{02WJ}{source}{}{sourcepointer-5}{}}
\newlabel{spaces-lemma-morphism-schemes-gives-representable-transformation-property@cref}{{[originalreference][0][]02WJ}{[1][0][]source}{}{}{}}
\makeatother
\crefname{originalreference}{original reference}{original references}
\Crefname{originalreference}{Original reference}{Original references}
\makeatletter
\newlabel{spaces-properties-lemma-subscheme}{{03JH}{source}{}{sourcepointer-6}{}}
\newlabel{spaces-properties-lemma-subscheme@cref}{{[originalreference][0][]03JH}{[1][0][]source}{}{}{}}
\makeatother
\title{The point conditions alpha through theta are Zariski local on the algebraic space}
\author{The Stacks Project Authors}
\date{stacks-project:decent-spaces.tex:local-snapshot-20261012; GNU FDL 1.2 or later}
\begin{document}
\maketitle
\noindent\textit{Editorial note.} The zeta case is a two-step definition unfolding carried as an explicit obligation; the theta case is a one-line citation.\par\medskip
\noindent\textit{Allowed standard prerequisites (editorial).}\begin{itemize}
\item Condition graph. lemma-bounded-fibres (carded L26703). Gives (gamma) = (alpha)+(beta).
\item Etale transfer for compositions. Spaces lemmas composition-representable-transformations-property, morphism-schemes-gives-representable-transformation-property and Morphisms lemmas composition-etale, open-immersion-etale (cross-chapter, named). The compositions U\_\{ij\} -> X stay etale in the epsilon case.
\item Open subschemes of schemes. Properties of Spaces, lemma subscheme (cross-chapter, named). The theta case.
\end{itemize}\medskip
\ifdefined\hypertarget\hypertarget{sourcepointer-1}{}\fi
\noindent\textit{Source reference.} Original lemma 03JX: sources/stacks-decent-spaces/pinned-original-source.tex, source lines 952--952. Navigation pointer only: the counted proof cites this statement by \textbackslash{}ref as a standing input; the deduction the counted proof has to make is carried out inside the counted span, and no part of that deduction is delegated to this pointer.\par\medskip
\ifdefined\hypertarget\hypertarget{sourcepointer-2}{}\fi
\noindent\textit{Source reference.} Original reference 02GN: focused-seeds/geometry/S9-decent-properties-local/ref-tags.txt, source lines 5--5. Navigation pointer only: the counted proof cites this statement by \textbackslash{}ref as a standing input; the deduction the counted proof has to make is carried out inside the counted span, and no part of that deduction is delegated to this pointer.\par\medskip
\ifdefined\hypertarget\hypertarget{sourcepointer-3}{}\fi
\noindent\textit{Source reference.} Original reference 02GP: focused-seeds/geometry/S9-decent-properties-local/ref-tags.txt, source lines 6--6. Navigation pointer only: the counted proof cites this statement by \textbackslash{}ref as a standing input; the deduction the counted proof has to make is carried out inside the counted span, and no part of that deduction is delegated to this pointer.\par\medskip
\ifdefined\hypertarget\hypertarget{sourcepointer-4}{}\fi
\noindent\textit{Source reference.} Original reference 02WK: focused-seeds/geometry/S9-decent-properties-local/ref-tags.txt, source lines 7--7. Navigation pointer only: the counted proof cites this statement by \textbackslash{}ref as a standing input; the deduction the counted proof has to make is carried out inside the counted span, and no part of that deduction is delegated to this pointer.\par\medskip
\ifdefined\hypertarget\hypertarget{sourcepointer-5}{}\fi
\noindent\textit{Source reference.} Original reference 02WJ: focused-seeds/geometry/S9-decent-properties-local/ref-tags.txt, source lines 8--8. Navigation pointer only: the counted proof cites this statement by \textbackslash{}ref as a standing input; the deduction the counted proof has to make is carried out inside the counted span, and no part of that deduction is delegated to this pointer.\par\medskip
\ifdefined\hypertarget\hypertarget{sourcepointer-6}{}\fi
\noindent\textit{Source reference.} Original reference 03JH: focused-seeds/geometry/S9-decent-properties-local/ref-tags.txt, source lines 9--9. Navigation pointer only: the counted proof cites this statement by \textbackslash{}ref as a standing input; the deduction the counted proof has to make is carried out inside the counted span, and no part of that deduction is delegated to this pointer.\par\medskip
\begin{lemma}
\label{lemma-properties-local}
Let $S$ be a scheme.
Let $\mathcal{P}$ be one of the properties
$(\alpha)$, $(\beta)$, $(\gamma)$, $(\delta)$, $(\epsilon)$, $(\zeta)$, or
$(\theta)$ of algebraic spaces listed in
Lemma \ref{lemma-bounded-fibres}.
Then if $X$ is an algebraic space over $S$, and $X = \bigcup X_i$ is a
Zariski open covering such that each $X_i$ has $\mathcal{P}$,
then $X$ has $\mathcal{P}$.
\end{lemma}

\begin{proof}
Let $X$ be an algebraic space over $S$, and let $X = \bigcup X_i$ is a
Zariski open covering such that each $X_i$ has $\mathcal{P}$.

\medskip\noindent
The case $\mathcal{P} = (\alpha)$. The condition $(\alpha)$ for $X_i$
means that for every $x \in |X_i|$ and every affine scheme $U$, and
\'etale morphism $\varphi : U \to X_i$ the fibre of $\varphi : |U| \to |X_i|$
over $x$ is finite. Consider $x \in X$, an affine scheme $U$ and
an \'etale morphism $U \to X$. Since $X = \bigcup X_i$ is a
Zariski open covering there exits a finite affine open covering
$U = U_1 \cup \ldots \cup U_n$ such that each $U_j \to X$ factors through
some $X_{i_j}$. By assumption the fibres of $|U_j | \to |X_{i_j}|$
over $x$ are finite for $j = 1, \ldots, n$. Clearly this means that
the fibre of $|U| \to |X|$ over $x$ is finite.
This proves the result for $(\alpha)$.

\medskip\noindent
The case $\mathcal{P} = (\beta)$. The condition $(\beta)$ for $X_i$ means
that every $x \in |X_i|$ is represented by a monomorphism from the
spectrum of a field towards $X_i$. Hence the same follows for $X$
as $X_i \to X$ is a monomorphism and $X = \bigcup X_i$.

\medskip\noindent
The case $\mathcal{P} = (\gamma)$.
Note that $(\gamma) = (\alpha) + (\beta)$ by
Lemma \ref{lemma-bounded-fibres}
hence the lemma for $(\gamma)$ follows from the cases treated above.

\medskip\noindent
The case $\mathcal{P} = (\delta)$. The condition $(\delta)$ for $X_i$ means
there exist schemes $U_{ij}$ and \'etale morphisms $U_{ij} \to X_i$ with
universally bounded fibres which cover $X_i$. These schemes also give an
\'etale surjective morphism $\coprod U_{ij} \to X$ and $U_{ij} \to X$
still has universally bounded fibres.

\medskip\noindent
The case $\mathcal{P} = (\epsilon)$. The condition $(\epsilon)$ for $X_i$ means
we can find a set $J_i$ and morphisms
$\varphi_{ij} : U_{ij} \to X_i$ such that each $\varphi_{ij}$
is \'etale, both projections $U_{ij} \times_{X_i} U_{ij} \to U_{ij}$
are quasi-compact, and $\coprod_{j \in J_i} U_{ij} \to X_i$ is surjective.
In this case the compositions $U_{ij} \to X_i \to X$ are \'etale
(combine
Morphisms, Lemmas
\ref{morphisms-lemma-composition-etale} and
\ref{morphisms-lemma-open-immersion-etale}
and
Spaces, Lemmas
\ref{spaces-lemma-composition-representable-transformations-property} and
\ref{spaces-lemma-morphism-schemes-gives-representable-transformation-property}
).
Since $X_i \subset X$ is a subspace we see that
$U_{ij} \times_{X_i} U_{ij} = U_{ij} \times_X U_{ij}$, and hence the
condition on fibre products is preserved. And clearly
$\coprod_{i, j} U_{ij} \to X$ is surjective. Hence $X$
satisfies $(\epsilon)$.

\medskip\noindent
The case $\mathcal{P} = (\zeta)$. The condition $(\zeta)$ for $X_i$
means that $X_i$ is Zariski locally quasi-separated. It is immediately
clear that this means $X$ is Zariski locally quasi-separated.

\medskip\noindent
For $(\theta)$, see
Properties of Spaces,
Lemma \ref{spaces-properties-lemma-subscheme}.
\end{proof}

\end{document}
